%%%% ABSTRACT
%%
%% Versão do resumo para idioma de divulgação internacional.

\begin{abstractutfpr}%% Ambiente abstractutfpr
With the increasing complexity of software systems and the need for fast and reliable deliveries, the DevOps approach has emerged to integrate development and operations teams, enhancing automation, collaboration, and software delivery. Within this approach, Continuous Integration (CI) and Continuous Delivery (CD) stand out as fundamental practices. Continuous Integration refers to the process of automating the integration of code from multiple developers into a shared repository, ensuring that frequent changes are continuously verified and tested. Continuous Delivery complements this process by automating software deployment to testing and production environments, enabling fast and reliable releases. To support these practices, CI/CD pipelines (also known as build and deployment pipelines) are used. These pipelines consist of automated workflows that systematically compile, test, and deploy applications, ensuring consistency and reproducibility throughout the software development lifecycle.
This study presents a comparative analysis of three leading CI/CD tools: GitHub Actions, GitLab CI/CD, and Jenkins. The methodology involves implementing CI/CD pipelines on each platform, followed by an assessment based on predefined criteria. To ensure result consistency and reliability, experiments will be conducted in a controlled environment where infrastructure variables remain constant. This research aims to provide insights into the capabilities of each tool, helping development teams choose the most suitable solution for automating their workflows.
\end{abstractutfpr}
